\chapter{Numerical and Floating-Point Semantics}
\label{app:numerics}

\chapterorient{Every algorithm in this book has been presented in exact
arithmetic---inner products are sums, the residual is $\vb-\mat{A}\vx$, a
polynomial is its value. A real machine does none of this exactly. It rounds, and
the order in which it rounds is visible in the last bits of the answer. This
appendix is about which of those rounding choices the reader may safely
\emph{ignore} and which are \emph{part of the algorithm}. The distinction is the
one \cref{ch:bigidea} drew between performance tricks and base algebra, sharpened
to the floating-point level: a \emph{transparent} trick computes the same result
faster and is invisible to the mathematics; a \emph{load-bearing} trick changes a
property the algorithm depends on---determinism, conditioning, IEEE
compliance---and must be stated as an explicit claim. We first give a working
floating-point primer (rounding, machine epsilon, non-associativity, cancellation,
conditioning), then catalogue the transparent tricks the book has quietly skipped
over, then treat the load-bearing ones in depth: non-associative reduction
orderings, fast-math reassociation, mixed precision, the Chebyshev
recurrence-versus-monomial choice (\cref{ch:smoothers}), and modified-versus-classical
Gram--Schmidt (\cref{ch:orthog}). \Cref{app:laws} catalogues which algebraic
\emph{laws} survive the rotation to pure functions; this appendix explains why
several of them hold only \emph{up to rounding}, and why that ``up to'' is itself a
specification.}

\section{A floating-point primer}
\label{sec:primer}

The reader who has written a numerical program has met the surprises of floating
point---a sum that depends on the order of its terms, a subtraction that loses all
its digits---without necessarily having a vocabulary for them. This section
supplies the vocabulary, at the level needed to follow the rest of the appendix.
Nothing here is deep; all of it is load-bearing for what follows.

\subsection{Floating-point numbers and rounding}

A binary floating-point number is a sign, a fraction (the \emph{significand} or
\emph{mantissa}), and an exponent: $x = \pm\, (1.b_1 b_2 \cdots b_{p-1})_2 \times
2^{e}$, with $p$ significand bits. IEEE-754 double precision (``binary64'',
the workhorse of every solver in this book) has $p = 53$ significand bits and an
11-bit exponent; single precision (``binary32'') has $p = 24$ and an 8-bit
exponent. The finite significand is the entire story: only numbers expressible in
$p$ bits exist, and every arithmetic result must be \emph{rounded} to the nearest
representable one.

\begin{definition}[Unit roundoff]
\label{def:eps}
The \emph{unit roundoff} (or \emph{machine epsilon}) $u$ is the largest relative
error incurred by rounding a real number to the nearest floating-point number:
under round-to-nearest, $\mathrm{fl}(x) = x(1+\delta)$ with $\abs{\delta} \le u$,
where $u = 2^{-p}$. For double precision $u = 2^{-53} \approx 1.11\times10^{-16}$;
for single precision $u = 2^{-24} \approx 5.96\times10^{-8}$.
\end{definition}

The fundamental model of floating-point arithmetic follows: each elementary
operation $\circ \in \{+,-,\times,\div\}$ returns the exactly-rounded result,
\begin{equation}
\label{eq:fpmodel}
  \mathrm{fl}(a \circ b) \;=\; (a \circ b)(1 + \delta), \qquad \abs{\delta} \le u .
\end{equation}
Each \emph{single} operation is as accurate as the format allows---a relative
error of at most $u$. The trouble is never one operation; it is the way errors
from many operations combine, and \eqref{eq:fpmodel} is the seed from which every
phenomenon below grows.

\begin{notationbox}
We write $\mathrm{fl}(\cdot)$ for ``the floating-point value of,'' $u$ for the unit
roundoff (\cref{def:eps}), and $\kappa$ for a condition number (\cref{def:cond}
below, and \cref{def:kappa} in \cref{ch:cg} for the matrix case). A relative error
``of order $u$'' means $O(u)$; ``of order $\kappa u$'' means the rounding has been
amplified by the conditioning. These three symbols---$u$, $\kappa$, and their
product $\kappa u$---carry the whole quantitative story of the appendix.
\end{notationbox}

\subsection{Addition is not associative}
\label{sec:nonassoc}

The single fact that drives most of this appendix is that floating-point addition,
though commutative, is \emph{not associative}:
\begin{equation}
\label{eq:nonassoc}
  \mathrm{fl}\bigl(\mathrm{fl}(a + b) + c\bigr)
  \;\ne\;
  \mathrm{fl}\bigl(a + \mathrm{fl}(b + c)\bigr)
  \quad\text{in general.}
\end{equation}
Each inner sum is rounded before the outer sum sees it, and the two roundings
discard different bits. A three-number example makes it concrete. In a toy decimal
format carrying three significant digits, take $a = 1.00$, $b = c = 5.00\times
10^{-3}$. Left-associating, $\mathrm{fl}(a+b) = \mathrm{fl}(1.005) = 1.00$ (the
$0.005$ falls off the third digit), and then $\mathrm{fl}(1.00 + 0.005) = 1.00$:
both small terms are lost. Right-associating, $\mathrm{fl}(b+c) =
\mathrm{fl}(0.010) = 0.010$ (no rounding---it fits), and then $\mathrm{fl}(1.00 +
0.010) = 1.01$: the two small terms, summed \emph{first}, survive their encounter
with the large one. Same three numbers, two bracketings, different answers in the
last digit. \Cref{fig:nonassoc} draws the two summation trees and the bits each
one keeps.

\begin{figure}[t]
\centering
\begin{tikzpicture}[font=\footnotesize]
  % ===== LEFT: left-leaning tree =====
  \begin{scope}
    \node[font=\small\bfseries, text=simBlue] at (1.3,3.1) {left-associated $((a{+}b){+}c)$};
    \node[data] (a) at (0,2.2) {$a$};
    \node[data] (b) at (1.0,2.2) {$b$};
    \node[data] (c) at (2.6,2.2) {$c$};
    \node[opbox, minimum width=9mm] (s1) at (0.5,1.2) {$+$};
    \node[opbox, minimum width=9mm] (s2) at (1.55,0.2) {$+$};
    \node[product, minimum width=14mm] (r1) at (1.55,-0.9) {sum$_L$};
    \draw[flow] (a) -- (s1); \draw[flow] (b) -- (s1);
    \draw[flow] (s1) -- (s2); \draw[flow] (c) -- (s2);
    \draw[flow] (s2) -- (r1);
    \node[font=\scriptsize, text=simRust, align=center] at (1.55,-1.7)
      {$a{+}b$ rounds first:\\small bits of $b$ lost early};
  \end{scope}
  \draw[simGray!50, thick] (3.8,-2.0) -- (3.8,3.3);
  % ===== RIGHT: right-leaning tree =====
  \begin{scope}[xshift=5.2cm]
    \node[font=\small\bfseries, text=simGreen] at (1.5,3.1) {right-associated $(a{+}(b{+}c))$};
    \node[data] (a2) at (0,2.2) {$a$};
    \node[data] (b2) at (1.6,2.2) {$b$};
    \node[data] (c2) at (2.6,2.2) {$c$};
    \node[opbox, minimum width=9mm] (t1) at (2.1,1.2) {$+$};
    \node[opbox, minimum width=9mm] (t2) at (1.05,0.2) {$+$};
    \node[product, minimum width=14mm] (r2) at (1.05,-0.9) {sum$_R$};
    \draw[flow] (b2) -- (t1); \draw[flow] (c2) -- (t1);
    \draw[flow] (a2) -- (t2); \draw[flow] (t1) -- (t2);
    \draw[flow] (t2) -- (r2);
    \node[font=\scriptsize, text=simGreen, align=center] at (1.05,-1.7)
      {$b{+}c$ summed first:\\small terms survive together};
  \end{scope}
  % the verdict band across the bottom
  \node[font=\scriptsize\itshape, text=simInk, align=center] at (3.8,-2.6)
    {$a=1.00,\ b=c=5{\times}10^{-3}$ (3-digit decimal):\quad
     sum$_L = 1.00 \ne 1.01 = $ sum$_R$};
\end{tikzpicture}
\caption[Two summation trees, two answers]{The non-associativity of
\eqref{eq:nonassoc} made visible. The same three numbers are summed by two
\emph{reduction trees}. \emph{Left:} left-association adds $a+b$ first; in a coarse
format the small term $b$ is swallowed by the large $a$ before $c$ is ever added,
and its contribution is lost. \emph{Right:} right-association adds the two small
terms $b+c$ first, so they accumulate to a magnitude that survives the encounter
with $a$. The two trees compute the identical mathematical sum $a+b+c$ but disagree
in the last digit, because rounding happens \emph{between} the additions and each
tree rounds at different places. A reduction is a tree, and the tree is observable
in the result.}
\label{fig:nonassoc}
\end{figure}

This is not pathology; it is the rule. A sum of $N$ numbers can be evaluated by any
of exponentially many trees---sequential left-to-right, pairwise (a balanced
binary tree), blocked, or whatever a parallel reduction happens to produce---and
the trees agree only in exact arithmetic. In floating point they agree to within a
bound of order $N u$ (sequential) or $u\log_2 N$ (pairwise) times the sum of
magnitudes, but they do not agree \emph{bit for bit}. \Cref{sec:reduction} returns
to this with the global reductions of CG and GMRES, where it becomes the central
load-bearing fact.

\subsection{Catastrophic cancellation}
\label{sec:cancel}

The second hazard is \emph{subtraction of nearly equal quantities}. When $a$ and
$b$ agree in their leading digits, the difference $a - b$ \emph{cancels} those
leading digits, and the result is built entirely from the low-order bits---the very
bits already corrupted by the roundings that produced $a$ and $b$. The relative
error in $a - b$ is magnified by roughly $\abs{a}/\abs{a-b}$, which is enormous
when the operands are close. \Cref{fig:cancel} shows the bookkeeping.

\begin{figure}[t]
\centering
\begin{tikzpicture}[font=\footnotesize, x=0.34cm, y=0.5cm]
  % a row of "digit cells": green = correct, gray = garbage
  \def\cell#1#2#3{\draw[fill=#3] (#1,#2) rectangle ++(0.9,0.8);}
  % a = 1.2345678  (all 7 digits correct, last is rounding-uncertain)
  \node[anchor=east, font=\scriptsize] at (-0.4,3.4) {$a = 1.2345678\underline{9}$};
  \foreach \i in {0,...,6}{\cell{\i*1}{3}{simBgGreen}}
  \cell{7}{3}{simBgGray}
  % b = 1.2345677  (agree in first 7)
  \node[anchor=east, font=\scriptsize] at (-0.4,2.2) {$b = 1.2345677\underline{2}$};
  \foreach \i in {0,...,6}{\cell{\i*1}{2}{simBgGreen}}
  \cell{7}{2}{simBgGray}
  % a-b = 0.0000001.7 : leading 7 cancel, only the last (garbage) survives, shifted up
  \node[anchor=east, font=\scriptsize] at (-0.4,0.9) {$a - b = 0.000000\underline{17}$};
  \foreach \i in {0,...,5}{\cell{\i*1}{0.6}{white}}
  \cell{6}{0.6}{simBgRust}
  \cell{7}{0.6}{simBgRust}
  % brace over the cancelled region
  \draw[decorate, decoration={brace, amplitude=4pt}, simGray]
    (6.9,3.95) -- (-0.05,3.95);
  \node[font=\scriptsize, text=simGray] at (3.4,4.45) {7 leading digits cancel};
  \draw[decorate, decoration={brace, amplitude=4pt, mirror}, simRust]
    (5.95,0.45) -- (7.85,0.45);
  \node[font=\scriptsize, text=simRust, align=center] at (9.6,0.9)
    {result built only from\\the \emph{uncertain} trailing bits};
  % legend
  \node[anchor=west, font=\scriptsize] at (0,-1.1)
    {\tikz\fill[simBgGreen](0,0)rectangle(0.3,0.3); correct\quad
     \tikz\fill[simBgGray](0,0)rectangle(0.3,0.3); rounding-uncertain\quad
     \tikz\fill[simBgRust](0,0)rectangle(0.3,0.3); amplified garbage};
\end{tikzpicture}
\caption[Catastrophic cancellation]{Subtracting two nearly equal numbers destroys
significant digits. The operands $a$ and $b$ agree in their first seven digits
(green); only their trailing digit (gray) carries rounding uncertainty. Their
difference $a-b$ \emph{cancels} the seven agreeing digits exactly, leaving a result
whose significant figures are built entirely from the previously-uncertain trailing
bits (rust)---which are then re-normalized up to the leading position, magnifying
their relative error by roughly $\abs{a}/\abs{a-b}$. The subtraction itself is
exactly rounded (its own $\delta$ is below $u$); the damage was done \emph{earlier},
when $a$ and $b$ were rounded, and cancellation merely exposes it. The cure is
always algebraic---reformulate to avoid the near-equal subtraction---never a more
accurate subtractor.}
\label{fig:cancel}
\end{figure}

\begin{pitfall}
Cancellation is widely \emph{mis}-diagnosed. The subtraction $a-b$ is itself
exactly rounded---its relative error is below $u$ like any other operation. What is
catastrophic is that it \emph{reveals} errors already present in $a$ and $b$ from
earlier roundings. Replacing the subtraction with a ``more accurate'' one buys
nothing; the fix is always to reformulate the computation so the near-equal
subtraction never occurs---summing positive terms before subtracting, using a
stable recurrence instead of an explicit difference (the Chebyshev story of
\cref{sec:cheby}), or rationalizing $\sqrt{x+1}-\sqrt{x}$ into
$1/(\sqrt{x+1}+\sqrt{x})$. The monomial form of a polynomial
(\cref{sec:cheby}) is catastrophic cancellation industrialized: large
alternating-sign terms that should sum to something small.
\end{pitfall}

\subsection{Conditioning: the problem versus the algorithm}
\label{sec:conditioning}

The last piece of vocabulary separates two sources of error that are constantly
conflated. The \emph{condition number} of a problem measures how much its
\emph{output} can change in response to a small change in its \emph{input}---a
property of the problem itself, before any algorithm is chosen.

\begin{definition}[Condition number of a problem]
\label{def:cond}
For a problem $y = f(x)$, the (relative) condition number is the smallest $\kappa$
such that a relative perturbation $\abs{\Delta x}/\abs{x}$ produces a relative
output change $\abs{\Delta y}/\abs{y} \lesssim \kappa\,\abs{\Delta x}/\abs{x}$. For
solving $\mat{A}\vx = \vb$ the relevant condition number is
$\kappa(\mat{A}) = \norm{\mat{A}}\,\norm{\mat{A}^{-1}}$, which for SPD $\mat{A}$ is
$\lambda_{\max}/\lambda_{\min}$ (\cref{def:kappa}).
\end{definition}

An algorithm is \emph{stable} (more precisely, \emph{backward stable}) if the
answer it computes is the exact answer to a slightly perturbed input---a
perturbation of relative size $O(u)$. Stability and conditioning then multiply:
\begin{equation}
\label{eq:errorrule}
  \underbrace{\frac{\norm{\hat{\vx} - \vx}}{\norm{\vx}}}_{\text{forward error}}
  \;\lesssim\;
  \underbrace{\kappa(\mat{A})}_{\text{conditioning of problem}}
  \times
  \underbrace{O(u)}_{\text{stability of algorithm}} .
\end{equation}
This single inequality is the organizing rule of numerical computing. A stable
algorithm on an ill-conditioned problem still returns a large forward error---not
because the algorithm misbehaved but because the problem amplifies any
perturbation, including the unavoidable $O(u)$ rounding of the input. An unstable
algorithm wastes accuracy even on a well-conditioned problem. The two failure modes
are independent, and \eqref{eq:errorrule} tells which is which: if $\kappa$ is huge,
no algorithm will save you (precondition the problem---\cref{ch:cg}); if $\kappa$ is
modest but the answer is still wrong, the algorithm is unstable (fix the algorithm,
as MGS fixes CGS in \cref{sec:orthog}).

\begin{keyidea}
Conditioning is a property of the \emph{problem}; stability is a property of the
\emph{algorithm}. Forward error $\lesssim \kappa \cdot O(u)$: the conditioning of
the problem times the (in)stability of the algorithm. The ill-conditioned operators
Palace assembles (fine meshes, high material contrast) make $\kappa$ large, which
is exactly why preconditioning (\cref{ch:cg}) and stable orthogonalization
(\cref{ch:orthog}) earn their keep---they attack the two factors separately.
\end{keyidea}

\section{Transparent tricks: invisible to the mathematics}
\label{sec:transparent}

With the primer in hand we can state the appendix's central distinction precisely
and then dispatch the easy half of it. A \emph{transparent} performance trick is a
rewrite that computes the \emph{same result}---bit for bit, or up to a benign
reordering whose error is of the same harmless order as the computation
already carries---but faster. The book has presented the unfolded, mathematically
transparent form throughout and relegated each such trick to a one-line note,
because by definition it changes nothing the mathematics can see.

\begin{definition}[Transparent performance trick]
\label{def:transparent}
A code transformation is \emph{transparent} if it is algebraically equivalent to
the unfolded form it replaces and does not change any property the algorithm
depends on---not the value (up to benign reordering), not the conditioning, not the
rounding behavior in any consumed way. Its only effect is speed or memory traffic.
A transparent trick may be \emph{ignored in the mathematical presentation} and
recorded as a footnote on the operation it accelerates.
\end{definition}

The catalogue of transparent tricks Palace uses, each algebraically equivalent to
its unfolded form:

\begin{itemize}[leftmargin=1.4em]
\item \textbf{Operator fusion.} Two consecutive elementwise passes (say a scale
  then an add) are fused into one loop that reads the input once and writes the
  output once, halving memory traffic. The arithmetic performed is identical; only
  the number of trips to memory changes. The \code{axpy} family of \cref{ch:cg} is
  a fused scale-and-add presented as one operator for exactly this reason.

\item \textbf{Tiling and blocking.} A large matrix operation is partitioned into
  cache-sized tiles processed one at a time, so the working set fits in fast memory.
  The set of multiply-adds is unchanged; only their \emph{schedule} changes. (When
  the schedule changes the \emph{reduction order} of a sum, the trick is no longer
  purely transparent---see \cref{sec:reduction}; blocking that keeps each reduction
  intact is transparent, blocking that re-trees a reduction is load-bearing.)

\item \textbf{Packing and batching.} Many small elementwise or per-element
  operations are gathered into one large contiguous kernel launch (the
  element-batched assembly of \cref{ch:matrixfree}). The per-element arithmetic is
  untouched; batching only amortizes launch and indexing overhead.

\item \textbf{Memory layout.} Choosing array-of-structs versus struct-of-arrays,
  or a blocked sparse format, changes which bytes are adjacent but not which
  numbers are computed.

\item \textbf{Recomputation versus lookup.} A cheap quantity (a Jacobian
  determinant at a quadrature point) may be recomputed on the fly rather than
  stored and reloaded, trading arithmetic for memory bandwidth. The recomputed
  value is bit-identical to the stored one (it is the same formula), so the choice
  is transparent.

\item \textbf{Sum-factorization} (\cref{ch:matrixfree}). High-order matrix-free
  operator application replaces the dense $O(P\cdot L)$ basis contraction by a
  sequence of one-dimensional sweeps, $O(d\,P^{1/d}L)$, exploiting the
  tensor-product structure of the basis. As \cref{ch:matrixfree} establishes at
  length, the two compute \emph{bit-for-bit the same quadrature-point values}---the
  product $\phi_i(\xi)\phi_j(\eta)$ is simply contracted one axis at a time---so
  sum-factorization is the canonical transparent trick: decisive for performance,
  invisible to the algebra.
\end{itemize}

\begin{keyidea}
A transparent trick is a \emph{factoring of the schedule}, not of the algorithm.
Fusion, tiling, packing, layout, recomputation, and sum-factorization all leave the
computed numbers unchanged (up to benign reordering) and change only how fast or
with how much memory traffic they are produced. The mathematics of the book
therefore presents the unfolded form and notes the trick in one line. The reader
who wants to understand \emph{what is computed} may ignore every transparent trick;
the reader who wants to understand \emph{why it is fast enough to run} should read
\cref{ch:matrixfree}.
\end{keyidea}

There is one boundary case worth flagging explicitly, because it is the bridge to
the next section. Tiling, blocking, and batching are transparent \emph{only so long
as they preserve the reduction order of every sum}. The moment a tiling re-brackets
a global reduction---computes partial sums per tile and then combines them in a
different tree than the sequential sweep would---it has crossed from transparent to
load-bearing, because by \cref{sec:nonassoc} it changes the last bits of the sum.
The same code transformation can be transparent or load-bearing depending on
whether a reduction tree hides inside it. This is precisely why the classification
is a matter of \emph{judgment}, flagged to the human when unclear (\cref{ch:bigidea}),
rather than a mechanical property of the code shape.

\section{Load-bearing tricks: part of the algorithm}
\label{sec:loadbearing}

A \emph{load-bearing} numerical trick is one that changes a property the algorithm
depends on---and therefore cannot be ignored in the mathematics, because ignoring
it would mis-state what the algorithm computes. Where a transparent trick gets a
footnote, a load-bearing trick gets an \emph{explicit algebraic claim}, naming the
property it buys.

\begin{definition}[Load-bearing numerical trick]
\label{def:loadbearing}
A numerical choice is \emph{load-bearing} if changing it changes a property of the
result that some consumer depends on: its \emph{value} in a way that matters (the
exact iteration count, the last bits when reproducibility is required), its
\emph{conditioning} or attained accuracy, its \emph{determinism}, or its
\emph{IEEE compliance}. A load-bearing trick must be stated in the specification as
an explicit claim, together with the property it secures, and may not be silently
``optimized'' into an algebraically-equivalent-in-exact-arithmetic alternative.
\end{definition}

\Cref{fig:taxonomy} sets the two kinds side by side as a decision the spec makes for
every numerical choice, and the rest of the section works through the load-bearing
column case by case.

\begin{figure}[t]
\centering
\begin{tikzpicture}[font=\footnotesize]
  % the input question
  \node[stage, minimum width=46mm] (q) at (0,2.4)
    {A numerical choice (reorder, reassociate,\\reduce precision, expand a polynomial)};
  % left branch: transparent
  \node[opbox, fill=simBgGreen, draw=simGreen, text width=42mm, align=left] (T) at (-3.6,0)
    {\textbf{Transparent}\\[2pt]
     same result (up to benign reorder),\\just faster: fusion, tiling, packing,\\
     layout, recompute, sum-factorization};
  \node[product, fill=simBgGreen, draw=simGreen, text width=34mm] (Tout) at (-3.6,-2.3)
    {\textbf{ignore in the math};\\ one-line note on the operation};
  % right branch: load-bearing
  \node[opbox, fill=simBgRust, draw=simRust, text width=44mm, align=left] (L) at (3.6,0)
    {\textbf{Load-bearing}\\[2pt]
     changes value / conditioning /\\determinism / IEEE compliance:\\
     reduction order, fast-math, mixed\\precision, recurrence-vs-monomial, MGS-vs-CGS};
  \node[product, fill=simBgRust, draw=simRust, text width=36mm] (Lout) at (3.6,-2.3)
    {\textbf{state as an algebraic claim};\\ name the property it buys};
  % edges
  \draw[flow] (q.south) -- (T.north)
    node[midway, above left, font=\scriptsize, text=simGreen]
      {same answer?};
  \draw[flow] (q.south) -- (L.north)
    node[midway, above right, font=\scriptsize, text=simRust]
      {changes a property?};
  \draw[flow] (T) -- (Tout);
  \draw[flow] (L) -- (Lout);
  % unclear -> human
  \node[extern, text width=30mm] (H) at (0,-3.7)
    {\emph{unclear}: flag to the human (\cref{ch:bigidea})};
  \draw[refedge] (T.south) .. controls +(0,-0.9) and +(-0.7,0.3) .. (H.west);
  \draw[refedge] (L.south) .. controls +(0,-0.9) and +(0.7,0.3) .. (H.east);
\end{tikzpicture}
\caption[Transparent versus load-bearing: the spec's decision]{The classification
every numerical choice passes through. \emph{Left (green):} if the choice computes
the same result up to a benign reordering and changes no depended-on property, it is
\emph{transparent}---ignored in the mathematical presentation, recorded as a one-line
note on the operation it accelerates. \emph{Right (rust):} if it changes the value
in a consumed way, the conditioning, the determinism, or the IEEE compliance, it is
\emph{load-bearing}---stated as an explicit algebraic claim naming the property it
secures. The hazard is mis-routing a load-bearing trick down the transparent path,
which silently changes the algorithm; when the routing is unclear the spec flags it
to the human (\cref{ch:bigidea}) rather than guessing. The five load-bearing cases
named on the right are the subject of \crefrange{sec:reduction}{sec:orthog}.}
\label{fig:taxonomy}
\end{figure}

\subsection{Non-associative reduction orderings}
\label{sec:reduction}

The first and most pervasive load-bearing trick is the one \cref{sec:numerics} of
\cref{ch:cg} named and deferred to here: the \emph{order} in which a global
reduction is summed. Every Krylov coefficient---the CG step length $\alpha$, the
direction coefficient $\beta$, the GMRES least-squares residual, every convergence
test---is built from inner products $\inner{\vu}{\vv} = \sum_i u_i v_i$, and an
inner product is a global reduction over all $N$ entries. By the non-associativity
of \cref{sec:nonassoc}, the value of that reduction depends on the summation tree.

The summation tree is not a fixed property of the program; it is chosen by the
runtime. A sequential loop sums left-to-right; a multithreaded reduction splits the
range into per-thread partial sums and combines them; a GPU reduction uses a
tree of warps and blocks; a different thread count or a different device produces a
different tree. Each tree gives a bit-different but equally valid value of the inner
product. \Cref{fig:redtree} contrasts two trees over the same eight terms.

\begin{figure}[t]
\centering
\begin{tikzpicture}[font=\footnotesize, scale=0.95]
  % ===== LEFT: sequential (left-leaning) reduction =====
  \begin{scope}
    \node[font=\small\bfseries, text=simBlue] at (1.6,3.5) {sequential: one long chain};
    \foreach \i [count=\j from 0] in {0,1,2,3,4,5,6,7}{
      \node[data, minimum width=4mm, inner sep=1.5pt] (x\i) at (\j*0.62,2.7)
        {\scriptsize$\i$};
    }
    % left-leaning accumulator chain
    \node[opbox, minimum width=4mm, inner sep=1.5pt] (acc) at (0.0,1.0) {\scriptsize$+$};
    \foreach \i in {1,...,7}{
      \pgfmathsetmacro\xx{\i*0.45}
      \node[opbox, minimum width=4mm, inner sep=1.5pt] (a\i) at (\xx,{1.0-\i*0.22}) {\scriptsize$+$};
    }
    \draw[flow] (x0)--(acc); \draw[flow] (x1)--(acc);
    \draw[flow] (acc)--(a2);
    \foreach \i [evaluate=\i as \p using int(\i-1)] in {2,...,7}{
      \ifnum\i>2 \draw[flow] (a\p)--(a\i); \fi
      \draw[flow] (x\i)--(a\i);
    }
    \node[product, minimum width=12mm] (rL) at (3.15,-0.9) {sum$_{\text{seq}}$};
    \draw[flow] (a7)--(rL);
    \node[font=\scriptsize, text=simRust, align=center] at (1.6,-1.7)
      {error $\sim N u$;\\ depth $N$, not parallel};
  \end{scope}
  \draw[simGray!50, thick] (4.6,-2.0) -- (4.6,3.7);
  % ===== RIGHT: pairwise (balanced tree) reduction =====
  \begin{scope}[xshift=5.6cm]
    \node[font=\small\bfseries, text=simGreen] at (2.2,3.5) {pairwise: balanced tree};
    \foreach \i [count=\j from 0] in {0,1,2,3,4,5,6,7}{
      \node[data, minimum width=4mm, inner sep=1.5pt] (y\i) at (\j*0.62,2.7)
        {\scriptsize$\i$};
    }
    % level 1: 4 adds
    \foreach \k [count=\m from 0] in {0,2,4,6}{
      \node[opbox, minimum width=4mm, inner sep=1.5pt] (l1\m) at ({(\k+0.5)*0.62},1.7) {\scriptsize$+$};
    }
    \draw[flow] (y0)--(l10); \draw[flow] (y1)--(l10);
    \draw[flow] (y2)--(l11); \draw[flow] (y3)--(l11);
    \draw[flow] (y4)--(l12); \draw[flow] (y5)--(l12);
    \draw[flow] (y6)--(l13); \draw[flow] (y7)--(l13);
    % level 2: 2 adds
    \node[opbox, minimum width=4mm, inner sep=1.5pt] (l20) at (1.24,0.7) {\scriptsize$+$};
    \node[opbox, minimum width=4mm, inner sep=1.5pt] (l21) at (3.10,0.7) {\scriptsize$+$};
    \draw[flow] (l10)--(l20); \draw[flow] (l11)--(l20);
    \draw[flow] (l12)--(l21); \draw[flow] (l13)--(l21);
    % level 3: 1 add
    \node[product, minimum width=12mm] (rR) at (2.17,-0.9) {sum$_{\text{tree}}$};
    \draw[flow] (l20)--(rR); \draw[flow] (l21)--(rR);
    \node[font=\scriptsize, text=simGreen, align=center] at (2.2,-1.7)
      {error $\sim u\log_2 N$;\\ depth $\log_2 N$, parallel};
  \end{scope}
\end{tikzpicture}
\caption[Two reduction trees over the same data]{The same eight-term reduction
$\sum_{i=0}^{7} z_i$ computed by two trees. \emph{Left:} a sequential accumulator
chain---the natural left-to-right loop---has depth $N$ and a worst-case error bound
of order $N u$ times the sum of magnitudes; it cannot be parallelized. \emph{Right:}
a balanced pairwise tree has depth $\log_2 N$, a tighter error bound of order
$u\log_2 N$, and maps onto parallel hardware. Both are correct reductions of the
same data; they round at different places and so generally produce \emph{different
last bits}. A multithreaded or GPU run picks a tree determined by its thread and
block geometry, so ``the same'' inner product computed on two devices---or with two
thread counts---returns bit-different values, and the iterative solver downstream
takes a bit-different (but equally valid) trajectory.}
\label{fig:redtree}
\end{figure}

The consequence for an iterative solver is precise and, at first, unsettling. Two
runs of the ``same'' CG solve on two machines---or on one machine with two thread
counts---compute bit-different inner products, hence bit-different $\alpha$ and
$\beta$, hence iterates differing in their last bits, and may even cross the
convergence tolerance at \emph{different iteration numbers}. The converged answer is
correct to tolerance in both cases; what is \emph{not} reproducible bit-for-bit is
the trajectory and the exact step count. This is the canonical load-bearing trick:
not a bug, but a property of the computation that a consumer may or may not depend
on, and that the specification must therefore \emph{name}.

\begin{keyidea}
An inner product is a reduction; a reduction is a tree; the tree determines the last
bits. Reduction order is \emph{load-bearing} precisely because it is observable: it
changes the bit-level trajectory and the exact iteration count of every Krylov
solver. The mathematics is identical across trees---each is a correct sum---so the
spec does not pick a tree, but it must \emph{state} that the trajectory is
tree-dependent, so that a downstream consumer knows reproducibility is not free.
\end{keyidea}

\begin{pitfall}
\textbf{Bit-reproducibility is not free, and is not on by default.} A consumer that
demands the identical bit pattern from run to run---a regression test asserting on
the last digits, a certification requiring reproducible output---must \emph{pin the
reduction order}: fix the thread count, fix the device, use a deterministic
(fixed-tree or compensated) reduction. Each of these costs performance: a
deterministic reduction forgoes the fastest tree the hardware offers, and a
fixed thread count forgoes load balancing. CG does not, and cannot, promise
bit-reproducibility for free---the same property that lets it run fast on any device
geometry is the property that makes its last bits device-dependent. The trade is
explicit: \emph{determinism costs throughput}, and the spec records reduction order
as the place that bill comes due.
\end{pitfall}

\subsection{Fast-math and reassociation flags}
\label{sec:fastmath}

A compiler in its default mode honors \eqref{eq:fpmodel} operation by operation and
preserves the parenthesization the source code wrote---it will not turn
$(a+b)+c$ into $a+(b+c)$, because by \cref{sec:nonassoc} that would change the
result. The various \emph{fast-math} flags (\code{-ffast-math},
\code{-funsafe-math-optimizations}, and their relatives) lift this restriction:
they license the compiler to treat floating-point arithmetic as if it were
associative, distributive, and exact, enabling reassociation, vectorization across
reduction boundaries, contraction of $a\cdot b + c$ into a fused multiply-add,
flushing denormals to zero, and assuming no NaNs or infinities appear.

\begin{table}[t]
\centering
\caption[What fast-math buys and breaks]{The fast-math bargain. Each license the
flag grants buys a performance opportunity and breaks a guarantee. The flags are
\emph{load-bearing}: enabling them changes the computed last bits and can change
which branch a NaN-guard takes, so they are a specification decision, not a free
optimization. Palace and the solvers of this book are compiled \emph{without}
unsafe fast-math precisely because the reduction-order and recurrence-stability
properties of \crefrange{sec:reduction}{sec:orthog} depend on the standard model
holding.}
\label{tab:fastmath}
\small
\setlength{\tabcolsep}{6pt}
\renewcommand{\arraystretch}{1.25}
\begin{tabular}{@{}p{0.30\linewidth}p{0.30\linewidth}p{0.32\linewidth}@{}}
\toprule
License granted & Buys & Breaks \\
\midrule
Reassociation ($+,\times$ associative) & vectorized / tree reductions, loop
  fusion & bit-reproducibility; a pinned reduction order (\cref{sec:reduction}) \\
Fused multiply-add contraction & one rounding instead of two; throughput &
  the exact $\mathrm{fl}(a\cdot b)$ a stable formula may rely on \\
Flush denormals to zero & no slow-path denormal handling & gradual underflow; tiny
  residuals snap to $0$ \\
Assume no NaN / Inf & cheaper comparisons, no guard code & NaN-propagation
  diagnostics; divergence detection \\
\bottomrule
\end{tabular}
\end{table}

\Cref{tab:fastmath} lays out the trade. The point for the specification is that
fast-math reassociation is \emph{exactly} the reduction-order freedom of
\cref{sec:reduction}, granted to the compiler rather than the runtime: it lets the
compiler pick a summation tree the source did not write. That makes it load-bearing
for the same reason---it changes last bits and reproducibility---and additionally
hazardous because the ``assume no NaN/Inf'' license can silently delete the very
guard code that an algorithm uses to detect divergence or a singular operator. A
solver that relies on a NaN propagating to its convergence check (a crude but real
divergence detector) will, under fast-math, march past the divergence unnoticed.
The disposition in this book: the standard model \eqref{eq:fpmodel} is assumed to
hold, unsafe fast-math is off, and any place that genuinely wants a fused
multiply-add or a relaxed reduction states it explicitly as the load-bearing choice
it is.

\subsection{Mixed precision: single apply, double accumulate}
\label{sec:mixed}

The third load-bearing trick trades precision for speed deliberately and in a
controlled place. Many operator applies and elementwise kernels can be carried in
\emph{single} precision---halving memory traffic and, on GPUs, often doubling or
quadrupling arithmetic throughput---provided the \emph{accumulation} that follows is
carried in \emph{double}. The pattern is ``single-precision apply, double-precision
accumulate,'' and it is the workhorse of modern mixed-precision iterative solvers.

\begin{figure}[t]
\centering
\begin{tikzpicture}[font=\footnotesize, node distance=8mm and 13mm]
  % single-precision lane (top)
  \node[data, fill=simBgGray, minimum width=14mm] (vin) {$\vv$\\\scriptsize fp32};
  \node[opbox, fill=simBgRust, draw=simRust, right=of vin, minimum width=22mm]
    (apply) {operator apply\\$\mat{A}\vv$\\\scriptsize \textbf{fp32}, fast};
  \node[data, fill=simBgGray, minimum width=16mm, right=of apply] (prod)
    {products $u_i v_i$\\\scriptsize fp32};
  \draw[flow] (vin) -- (apply);
  \draw[flow] (apply) -- (prod);
  % the precision-promotion boundary
  \draw[densely dashed, simViolet, thick] (prod.north east)++(0.35,0.7)
    -- ($(prod.south east)+(0.35,-0.7)$);
  \node[font=\scriptsize, text=simViolet, anchor=south, rotate=90] at
    ($(prod.east)+(0.55,0)$) {promote fp32$\to$fp64};
  % double-precision accumulate (right)
  \node[opbox, fill=simBgBlue, draw=simBlue, right=18mm of prod, minimum width=24mm]
    (acc) {reduction\\$\sum_i u_i v_i$\\\scriptsize \textbf{fp64}, accurate};
  \node[product, right=of acc, minimum width=16mm] (out) {$\alpha,\beta,\norm{\vr}$\\\scriptsize fp64};
  \draw[flow] (prod) -- (acc);
  \draw[flow] (acc) -- (out);
  % annotations
  \node[font=\scriptsize, text=simRust, align=center, below=2mm of apply]
    {bulk work:\\cheap, low precision};
  \node[font=\scriptsize, text=simBlue, align=center, below=2mm of acc]
    {the reduction:\\full precision};
\end{tikzpicture}
\caption[Single-precision apply, double-precision accumulate]{The mixed-precision
dataflow. The expensive bulk work---the operator apply $\mat{A}\vv$ and the
elementwise products $u_i v_i$ (rust)---is carried in single precision (fp32),
where it is cheap in both memory bandwidth and throughput. At the
\emph{accumulation} boundary (violet, dashed) the products are promoted to double
precision (fp64), and the global reduction $\sum_i u_i v_i$ that forms each Krylov
coefficient is summed in double (blue). This is load-bearing, not transparent: the
result is \emph{not} the bit-equal double-precision answer, and the soundness of the
trade rests on the conditioning argument of \cref{sec:conditioning}---the
single-precision apply injects a relative error of order $u_{32}\sim 6\times10^{-8}$
that the conditioning of the problem amplifies, so the trick is safe only while
$\kappa\,u_{32}$ stays below the target tolerance.}
\label{fig:mixed}
\end{figure}

Why accumulate in double? Because the accumulation is the global reduction of
\cref{sec:reduction}, and a reduction of $N$ single-precision terms accumulates an
error of order $N u_{32}$---which, for $N$ in the millions and $u_{32}\approx
6\times10^{-8}$, overwhelms the answer. Promoting the running sum to double drops
the accumulation error to order $N u_{64} \approx N\cdot10^{-16}$, negligible again,
while the bulk arithmetic stays cheap. The single-precision \emph{apply} still
injects a relative error of order $u_{32}$ into each matrix-vector product, and here
the conditioning rule \eqref{eq:errorrule} governs: that error is amplified by
$\kappa(\mat{A})$, so the mixed-precision solve attains a forward accuracy of order
$\kappa\,u_{32}$ rather than $\kappa\,u_{64}$. The trade is sound \emph{exactly
when} $\kappa\,u_{32}$ is comfortably below the solve tolerance---a well-conditioned
or well-preconditioned system can run its applies in single and lose nothing
visible; an ill-conditioned one cannot.

\begin{keyidea}
Mixed precision is the conditioning rule \eqref{eq:errorrule} turned into an
engineering knob. Do the \emph{bulk} arithmetic (the apply, the products) in single
precision for speed, but carry the \emph{accumulation} (the reduction) in double so
the global sum does not lose its low bits. The result is load-bearing---it is not
the double-precision answer---and it is sound only while $\kappa\,u_{32}$ stays below
tolerance. This is why mixed-precision iterative refinement and single/double Krylov
solvers pair so naturally with preconditioning (\cref{ch:cg}): the preconditioner
shrinks $\kappa$, which is exactly what buys the headroom to drop the apply to
single precision.
\end{keyidea}

\subsection{The Chebyshev recurrence versus the monomial form}
\label{sec:cheby}

The fourth load-bearing trick is a case where two algebraically identical
evaluations of the \emph{same polynomial} differ by everything that matters
numerically. The Chebyshev smoother of \cref{ch:smoothers} applies a fixed-degree
polynomial $p_k(\mat{D}^{-1}\mat{A})$ to the residual. There are two ways to
evaluate it, and the choice between them is the algorithm.

The \emph{monomial form} writes the polynomial out in powers and sums them:
\begin{equation}
\label{eq:monomial}
  p_k(\mat{B})\,\vr \;=\; \sum_{j=0}^{k} c_j\,\mat{B}^{j}\vr,
  \qquad \mat{B} = \mat{D}^{-1}\mat{A},
\end{equation}
forming $\mat{B}\vr, \mat{B}^2\vr, \dots$ and accumulating with the coefficients
$c_j$. The \emph{three-term recurrence} (a Clenshaw/Horner-style evaluation) never
forms the powers or the coefficients; it threads three vectors through a short loop,
as \cref{ch:smoothers} derives. In exact arithmetic the two compute the identical
vector. In floating point they could not be more different.

The monomial form \eqref{eq:monomial} is numerically catastrophic for the degrees in
use. The high powers $\mat{B}^{j}\vr$ grow enormous (their size is governed by the
largest eigenvalue raised to the $j$), the coefficients $c_j$ alternate in sign and
also grow large, and the sum of large alternating-sign terms collapsing to a modest
result is catastrophic cancellation (\cref{sec:cancel}) at industrial scale: the
answer is built from the differences of huge numbers whose low bits are pure
rounding noise. The three-term recurrence sidesteps this entirely---it evaluates the
same polynomial through a sequence of moderate, well-scaled intermediate vectors,
never forming a large power or a large coefficient, so no cancellation occurs. This
is the stability that Phillips and Fischer analyze and that Palace follows
directly~\citep[\S2]{phillipsfischer2022}. \Cref{fig:cheby} plots the divergence.

\begin{figure}[t]
\centering
\begin{tikzpicture}
\begin{semilogyaxis}[
  width=12.0cm, height=6.6cm,
  xlabel={polynomial degree $k$},
  ylabel={relative error $\norm{\hat{p}_k(\mat{B})\vr - p_k(\mat{B})\vr}/\norm{p_k(\mat{B})\vr}$},
  ymode=log, ymin=1e-17, ymax=1e6,
  xmin=1, xmax=20,
  grid=both, grid style={simGray!18},
  legend pos=north west, legend cell align=left,
  legend style={font=\footnotesize, draw=simGray!40},
  label style={font=\footnotesize}, tick label style={font=\scriptsize},
  ytick={1e-16,1e-12,1e-8,1e-4,1e0,1e4},
  every axis plot/.append style={very thick}]
  % monomial: error blows up roughly like (cond)^k * u -> exponential growth
  \addplot[simRust, mark=*, mark size=1.2pt] coordinates {
    (1,2e-16)(2,1e-15)(3,8e-15)(4,9e-14)(5,1e-12)(6,2e-11)(7,3e-10)(8,5e-9)
    (9,8e-8)(10,1e-6)(11,2e-5)(12,3e-4)(13,5e-3)(14,8e-2)(15,1.0)(16,1e1)
    (17,2e2)(18,4e3)(19,7e4)(20,1e6)};
  \addlegendentry{monomial $\sum_j c_j \mat{B}^j\vr$ (unstable)}
  % recurrence: stays flat at machine precision (grows only mildly with k)
  \addplot[simGreen, mark=triangle*, mark size=1.4pt] coordinates {
    (1,2e-16)(2,2e-16)(3,3e-16)(4,3e-16)(5,4e-16)(6,4e-16)(7,5e-16)(8,5e-16)
    (9,6e-16)(10,6e-16)(11,7e-16)(12,7e-16)(13,8e-16)(14,8e-16)(15,9e-16)
    (16,9e-16)(17,1.0e-15)(18,1.1e-15)(19,1.2e-15)(20,1.3e-15)};
  \addlegendentry{three-term recurrence (stable)}
  % machine-epsilon reference + the "answer is garbage" line at 1.0
  \addplot[simGray, densely dashed, domain=1:20, samples=2] {2.2e-16};
  \node[font=\scriptsize, text=simGray, anchor=west] at (axis cs:1.2,5e-16) {machine $u$};
  \addplot[simInk, densely dotted, domain=1:20, samples=2] {1.0};
  \node[font=\scriptsize, text=simInk, anchor=west] at (axis cs:1.2,2.2) {no digits left};
\end{semilogyaxis}
\end{tikzpicture}
\caption[Chebyshev recurrence versus monomial evaluation]{Relative error of the two
evaluations of the same degree-$k$ Chebyshev polynomial, as the degree grows, on a
log scale (schematic of the Phillips--Fischer experiment). The \emph{monomial form}
(rust) accumulates error that grows roughly geometrically in $k$---the high powers
$\mat{B}^j\vr$ and the alternating coefficients $c_j$ produce catastrophic
cancellation (\cref{sec:cancel}), and by a modest degree the result has \emph{no
correct digits left} (the dotted line at relative error $1$). The \emph{three-term
recurrence} (green) evaluates the identical polynomial through well-scaled
intermediate vectors and stays pinned near machine precision $u$ for every degree.
The two forms are equal in exact arithmetic and not interchangeable in floating
point: the sequentiality of the recurrence is load-bearing, not a convenience, and
choosing the monomial form ``to vectorize the powers'' silently destroys the
smoother~\citep{phillipsfischer2022}.}
\label{fig:cheby}
\end{figure}

\begin{pitfall}
The recurrence form looks like a \emph{sequential bottleneck} begging to be
parallelized: each step reads the previous step's vectors. The temptation is to
``unroll'' it into the monomial form \eqref{eq:monomial}, where the powers
$\mat{B}^j\vr$ can be precomputed independently. \emph{Do not}---this is the same
trap as replacing MGS with CGS (\cref{sec:orthog}). The sequentiality of the
recurrence is not an obstacle to the algorithm; it \emph{is} the algorithm's
stability. Expanding to monomials trades a true sequential dependency for
catastrophic cancellation, computing the same polynomial in exact arithmetic and
garbage in floating point. The spec records the recurrence form as a load-bearing
algebraic claim (``evaluate $p_k$ by the three-term recurrence, not the monomial
expansion''), with the property it buys (floating-point stability) named
explicitly.
\end{pitfall}

\subsection{Modified versus classical Gram--Schmidt}
\label{sec:orthog}

The fifth load-bearing trick closes the loop with \cref{ch:orthog}, where it was
derived in full; here we file it in the taxonomy. Orthogonalizing a vector against a
basis can subtract the projections \emph{all at once against the original}
(classical Gram--Schmidt, CGS) or \emph{one at a time against the running residual}
(modified Gram--Schmidt, MGS). The two compute the identical orthogonal residual in
exact arithmetic---both are the projection $(\mat{I} - \mat{Q}\mat{Q}^{\!\top})\vw$.
In floating point they are not interchangeable: CGS loses orthogonality
\emph{quadratically} in the basis condition number, $\eta\sim\kappa^2 u$, while MGS
loses it only \emph{linearly}, $\eta\sim\kappa u$ (\cref{ch:orthog},
\cref{def:cond} here for $\kappa$).

The structural fact, drawn in full in \cref{ch:orthog}, is that CGS's independence
(all projections read the same input, so they fan out in parallel) is exactly what
makes it unstable, and MGS's sequential dependency (each projection reads the
freshly-cleaned residual) is exactly what makes it stable. The sequential dependency
is \emph{load-bearing}: it cannot be optimized into the parallel CGS form without
silently degrading the basis, and the degradation surfaces only later as a solver
that stalls or a spurious eigenvalue. The ``twice is enough'' repair (CGS2:
re-run the parallel classical pass a second time) recovers full orthogonality
$\eta\sim u$ at double the cost, recovering MGS-grade stability while keeping
batched, parallel-friendly reductions---the sweet spot of \cref{ch:orthog}'s
trade-off table. The cross-reference is the point: this is the same
transparent-versus-load-bearing decision, and MGS's column order lands firmly in the
load-bearing column, which is why \cref{ch:orthog} treats the order as a sequential
obstruction rather than an inefficiency to be removed.

\section{How the book has handled this throughout}
\label{sec:throughout}

The distinction of this appendix has been operating, quietly, in every chapter. The
governing rule was stated in \cref{ch:bigidea}: a significant fraction of Palace's
C++ exists because it was tuned for CPU cache and SIMD, a cost model that is not the
target's, and most of that code shape is a \emph{performance trick} to be separated
from the \emph{base algebra}. This appendix is the floating-point refinement of that
separation.

The mechanism throughout has been uniform. The pure-function form the book presents
(the ``L1 form,'' in the layered vocabulary) is always the \emph{unfolded} form---the
mathematics with every transparent trick removed. Sum-factorization is presented as a
dense contraction and footnoted (\cref{ch:matrixfree}); fusion is presented as
separate scale-and-add operators conceptually and fused only in the note; tiling and
layout never appear in the algebra at all. A reader who follows only the unfolded
forms computes the right answer, just slower---which is the definition of transparent
(\cref{def:transparent}).

The load-bearing tricks, by contrast, each earned an \emph{explicit note} at the
point of use, stating the property bought:
\begin{itemize}[leftmargin=1.4em, nosep]
\item CG's inner products carry the reduction-order note (\cref{ch:cg},
  \cref{sec:numerics}): the trajectory is tree-dependent; reproducibility must be
  bought.
\item The Chebyshev smoother carries the recurrence-versus-monomial note
  (\cref{ch:smoothers}): the three-term recurrence is the stable evaluation; the
  sequentiality is load-bearing.
\item The orthogonalizer carries the MGS-versus-CGS note (\cref{ch:orthog}): the
  column order is a sequential dependency, not an inefficiency.
\item The preconditioned and mixed-precision solvers carry the conditioning note
  (this appendix, \cref{sec:mixed}): the attainable accuracy is $\kappa$ times the
  apply precision.
\end{itemize}
That asymmetry---transparent tricks footnoted, load-bearing tricks claimed---is the
book's consistent discipline, and \cref{app:laws} relies on it: several of the
algebraic ``laws'' catalogued there (commutativity of a reduction, associativity of
an accumulation) hold \emph{only up to rounding}, and the reason a law is qualified
``up to reduction order'' rather than asserted outright is exactly the load-bearing
status of \cref{sec:reduction}. The non-law is the specification.

\section{Practical guidance: when to care}
\label{sec:guidance}

The appendix closes with the engineering question a reader actually faces: of all
this, when does it matter? The honest answer is that for most of a simulation it does
not, and knowing \emph{where} it does is the skill.

\begin{itemize}[leftmargin=1.4em]
\item \textbf{Iterative-solver convergence.} Reduction order, fast-math, and mixed
  precision all perturb the Krylov trajectory. The \emph{converged} answer is robust
  (it is correct to tolerance regardless), but the \emph{iteration count} and the
  \emph{last bits} are not. Care when the iteration count is itself a measured
  quantity, when a tolerance is set near machine precision (where the trajectory
  noise approaches the signal), or when a solve is near the edge of converging at
  all---there the choice of reduction order can decide convergence versus stagnation.

\item \textbf{Eigenvalue accuracy.} Orthogonalization quality is load-bearing for
  eigensolvers in a way it is not for linear solves. A linear solve tolerates a
  somewhat degraded Krylov basis (it re-converges); an eigensolver reads the basis
  directly, and a lost-orthogonality basis manufactures \emph{spurious eigenvalues}
  (\cref{ch:orthog}). Care: use CGS2 or re-orthogonalized MGS whenever eigenvalues
  are the product, never bare CGS.

\item \textbf{Reproducibility requirements.} If a downstream consumer asserts on
  exact bits---a regression test, a certification, a bit-exact replay---reduction
  order must be pinned and fast-math disabled, at a throughput cost
  (\cref{sec:reduction}). Care: decide reproducibility as a requirement \emph{up
  front}, because retrofitting determinism onto a solver tuned for the fastest tree
  is expensive.

\item \textbf{When \emph{not} to care.} For a single well-conditioned solve whose
  output is a field or a scalar reported to engineering precision, with no
  bit-exact reproducibility requirement, every transparent trick may be taken freely
  and the load-bearing ones tuned for speed: single-precision applies, the fastest
  reduction tree, aggressive fusion. The vast majority of a simulation lives here.
  The discipline is not paranoia everywhere; it is knowing the handful of places---near
  convergence, at eigenvalues, under reproducibility mandates---where the last bits
  stop being noise and start being the answer.
\end{itemize}

\begin{keyidea}
\textbf{The one rule to carry away.} A \emph{transparent} trick is invisible to the
mathematics---it computes the same result faster, so the book presents the unfolded
form and ignores the trick. A \emph{load-bearing} trick \emph{is} the
mathematics---it changes determinism, conditioning, or IEEE compliance, so the book
states it as an explicit algebraic claim naming the property it buys. Mis-classifying
a load-bearing trick as transparent (replacing MGS by CGS, the recurrence by
monomials, double accumulation by single, a pinned reduction by the fastest tree)
silently changes the algorithm. When in doubt, the choice is flagged to the human,
because the cost of guessing wrong is an answer that is plausibly, invisibly, wrong.
\end{keyidea}

\summarysection

A real machine rounds, and the order in which it rounds is observable in the last
bits. This appendix separated the rounding choices that the mathematics may
\emph{ignore} from those that \emph{are} the mathematics. The primer
(\cref{sec:primer}) established the working facts: the unit roundoff $u$
(\cref{def:eps}), the non-associativity of floating-point addition
(\cref{eq:nonassoc}), catastrophic cancellation (\cref{sec:cancel}), and the
conditioning rule, forward error $\lesssim \kappa\cdot O(u)$ (\cref{eq:errorrule}),
which separates a property of the \emph{problem} from a property of the
\emph{algorithm}. \emph{Transparent} tricks (\cref{sec:transparent})---fusion,
tiling, packing, layout, recomputation, sum-factorization---compute the same result
faster and are footnoted, not stated. \emph{Load-bearing} tricks
(\cref{sec:loadbearing}) change a depended-on property and are stated as claims:
non-associative reduction order makes the Krylov trajectory device-dependent and
bit-reproducibility a paid-for property (\cref{sec:reduction}); fast-math grants the
compiler the same reordering freedom and breaks reproducibility and NaN-guards
(\cref{sec:fastmath}); mixed precision does the bulk work in single and the
accumulation in double, sound only while $\kappa\,u_{32}$ stays below tolerance
(\cref{sec:mixed}); the Chebyshev three-term recurrence is the stable evaluation and
the monomial form is catastrophic cancellation (\cref{sec:cheby}); and the MGS
column order is a load-bearing sequential dependency, not an inefficiency
(\cref{sec:orthog}). The book has applied this discipline throughout
(\cref{sec:throughout})---unfolded forms for the transparent, explicit notes for the
load-bearing---and \cref{app:laws} depends on it, since several algebraic laws hold
only ``up to reduction order,'' a qualification that is itself a specification. The
practical lesson (\cref{sec:guidance}): care near convergence, at eigenvalues, and
under reproducibility mandates; elsewhere, take the speed.

\furtherreading

The standard reference for the analysis of every phenomenon in this appendix---the
floating-point model \eqref{eq:fpmodel}, error bounds for summation and inner
products, the stability of Gram--Schmidt and of polynomial evaluation---is Higham's
treatment as developed in the numerical-linear-algebra texts; the conditioning and
stability framing of \cref{sec:conditioning}, including the forward-error rule
\eqref{eq:errorrule} and the backward-stability viewpoint, follows
Trefethen and Bau~\citep{trefethenbau1997}, whose first lectures are the cleanest
short account of conditioning, stability, and floating point in print. The
summation-tree error bounds, the quantitative loss-of-orthogonality results for CGS
versus MGS ($\kappa^2 u$ versus $\kappa u$), the ``twice is enough''
re-orthogonalization theorem, and the matrix condition number are developed in full
in Golub and Van Loan~\citep{golubvanloan2013}, the definitive matrix-computations
reference and the source for the orthogonalization trade-off of \cref{ch:orthog}.
The load-bearing stability of the Chebyshev three-term recurrence over its monomial
expansion, the coefficient formulas, and the analysis Palace follows directly are due
to Phillips and Fischer~\citep{phillipsfischer2022}; their \S2 is the precise
statement that the sequentiality of the recurrence is part of the algorithm. For the
broader principle---separating performance tricks from base algebra, of which this
appendix is the floating-point instance---see \cref{ch:bigidea}; for the algebraic
laws that this appendix qualifies ``up to rounding,'' see \cref{app:laws}.
